\section*{Заключение}
\addcontentsline{toc}{section}{Заключение}

В работе разработан и экспериментально апробирован прототип системы,
интегрирующей профилирование производительности Go-приложений в
CI/CD-пайплайн для раннего обнаружения регрессий. Поставленная цель —
совместить статистическое сравнение агрегированных метрик бенчмарков
(\texttt{benchstat}) с функциональной диагностикой через pprof-профили в
едином, автоматически выполняемом на каждый pull request шаге CI —
достигнута: реализованы оба механизма (раздел~\ref{sec:razdel4}),
формализован критерий их совместного применения
(раздел~\ref{sec:razdel2}), а на четырёх управляемых сценариях
(раздел~\ref{sec:razdel5}) показано, что оба уровня данных дополняют друг
друга — сценарий \texttt{quadratic-dedup} продемонстрировал регрессию по
времени выполнения, сопровождающуюся \emph{улучшением} метрики аллокаций,
что делает диагностику по одному только агрегату или одному только
профилю недостаточной.

Все задачи, сформулированные в разделе~\ref{sec:razdel1}, решены: проведён
анализ инструментов и подходов, формализована модель процесса и критерий
обнаружения регрессии, спроектирована и реализована архитектура
прототипа, проведена апробация с оценкой корректности обнаружения и
накладных расходов профилирования (эмпирически $\approx$2,2\%,
статистически неотличимо от шума измерения).

Практическая значимость — в переносимости слоя сравнения
(\texttt{scripts/}), взаимодействующего с проверяемым кодом только через
стандартный формат вывода \texttt{go test -bench} и применимого к
произвольному Go-пакету без изменений. Теоретическая значимость — в
формализации критерия~\eqref{eq:delta}--\eqref{eq:fail}, разделяющего
статистическую и практическую значимость отклонения, и в разграничении
"непрерывного профилирования" на процессный (встроенный в CI) и
эксплуатационный (мониторинг production) смыслы — работа реализует
первый.

Направления дальнейших исследований: проведение экспериментов
раздела~\ref{sec:razdel5} в полном масштабе для точных доверительных
интервалов; отдельный механизм обнаружения изменения класса
алгоритмической сложности через бенчмарк с варьируемым размером входа, а
не фиксированным; применимость прототипа к конкурентному коду, требующему
анализа через \texttt{go tool trace}, не рассмотренного в данной работе.
